\subsection{Examples of formally smooth algebras}

Let k be a ring.

\begin{enumerate}
\def\labelenumi{\arabic{enumi})}
\tightlist
\item
  Let P be a projective \(k\)-module. The symmetric \(k\)-algebra \(\mathbf{S}_k(\mathrm{P})\) is formally smooth for the discrete topology, and a fortiori for the topology defined by its grading. Indeed, for any \(k\)-algebra C and every ideal N of C, the algebra homomorphisms from \(\mathbf{S}_k(\mathrm{P})\) into C (resp. C/N) are in bijective correspondence with the \(k\)-linear maps from P into C (resp. C/N), and the canonical map \(\mathrm{Hom}_k(\mathrm{P},\mathrm{C}) \to \mathrm{Hom}_k(\mathrm{P},\mathrm{C}/\mathrm{N})\) is surjective.
\end{enumerate}

Therefore (Proposition 3, c)), the \(k\)-algebra \(\widehat{\mathbf{S}}_k(\mathrm{P}) = \prod_{n\geqslant 0}\mathbf{S}_k^n (\mathrm{P})\) is formally smooth (for the product topology of the discrete topologies on the \(\mathbf{S}_k^n (\mathrm{P}))\)); indeed, it is the completion of the \(k\)-algebra \(\mathbf{S}_k(\mathrm{P})\) for the topology defined by the grading.

\begin{enumerate}
\def\labelenumi{\arabic{enumi})}
\setcounter{enumi}{1}
\item
  For any family of indeterminates \(\mathbf{T} = (\mathrm{T}_i)_{i\in I}\), the \(k\)-algebra of polynomials \(k[\mathbf{T}]\) and the \(k\)-algebra of formal power series \(k[[\mathbf{T}]]\), endowed with its canonical topology, are formally smooth; this follows from Example 1. If \(k\) is a field, the purely transcendental extension \(k(\mathbf{T})\) is formally smooth (No. 2, Proposition 4 a)).
\item
  Let \(f \in k[\mathrm{T}]\) be a polynomial in one indeterminate. To say that the \(k\)-algebra \(k[\mathrm{T}]/(f)\) is formally smooth is to say that the following property is satisfied: for every \(k\)-algebra C and every ideal N of square zero in C, every root of \(f\) in C/N lifts to a root of \(f\) in C. This is the case when \(f\) and its derivative \(f'\) generate the unit ideal. Indeed, let \(\alpha\) be a root of \(f\) in C/N and let \(a\) be an element of C lifting \(\alpha\). Then \(f(a)\) belongs to N and consequently \(f'(a)\) is invertible in C; the element \(b = a - f'(a)^{-1}f(a)\) lifts \(\alpha\). Since \(f'(a)^{-1}f(a)\) is of square zero, we have
\end{enumerate}

\[
f (b) = f (a) - f ^ {\prime} (a) f ^ {\prime} (a) ^ {- 1} f (a) = 0.
\]

THEOREM 1 (I. S. Cohen). Let \(k\) be a field and \(K\) a separable extension of \(k\). Then \(K\) is a formally smooth \(k\)-algebra.

Let C be a \(k\)-algebra, N an ideal of square zero in C, \(\pi: \mathrm{C} \to \mathrm{C/N}\) the canonical homomorphism, and \(\varphi: \mathrm{K} \to \mathrm{C/N}\) a homomorphism of \(k\)-algebras. It remains to construct a lifting of \(\varphi\). We distinguish two cases.

\begin{enumerate}
\def\labelenumi{\Alph{enumi})}
\item
  Suppose first that k has characteristic 0. Consider the pairs \((\mathrm{K}^{\prime}, \tilde{\varphi}^{\prime})\), where \(K^{\prime}\) is a subextension of K and \(\tilde{\varphi}^{\prime}: K^{\prime} \to C\) is a lifting of the restriction of \(\varphi\) to \(K^{\prime}\). The set of these pairs, ordered by extension, is inductive; by Zorn's lemma (E, III, p.~20, th. 2), there exists a maximal pair \((\mathrm{K}^{\prime}, \tilde{\varphi}^{\prime})\). Let us prove that \(K^{\prime}\) is equal to K. Let \(x \in K - K^{\prime}\). If x is transcendental over \(K^{\prime}\), the \(K^{\prime}\)-algebra \(\mathrm{K}^{\prime}(x)\) is formally smooth (example 2). If x is algebraic over \(K^{\prime}\), its minimal polynomial \(f \in K^{\prime}[T]\) is relatively prime to its derivative (A, V, p.~37, prop. 4), and \(\mathrm{K}^{\prime}(x)\) identifies with the \(K^{\prime}\)-algebra \(\mathrm{K}^{\prime}[\mathrm{T}]/(f)\), which is therefore a formally smooth \(K^{\prime}\)-algebra (example 3). In both cases, \(\mathrm{K}^{\prime}(x)\) is formally smooth over \(K^{\prime}\), and there exists an extension of \(\tilde{\varphi}^{\prime}\) to \(\mathrm{K}^{\prime}(x)\) that lifts the restriction of \(\varphi\) to \(\mathrm{K}^{\prime}(x)\), contradicting the maximality of \((\mathrm{K}^{\prime}, \tilde{\varphi}^{\prime})\).
\item
  Suppose \(k\) has characteristic \(p \neq 0\). Consider the ring homomorphism F: C → C such that F(x) = x \(^{p}\); one has F(x) = 0 for x ∈ N, so that there exists a unique ring homomorphism λ: C/N → C such that λ ∘ π = F. One has π(λ(π(x))) = π(x \(^{p}\)) = π(x) \(^{p}\); since π is surjective, one therefore has π(λ(z)) = z \(^{p}\) for every element z of C/N. On the other hand, let f: K → K \(^{p}\) be the isomorphism y ↦ y \(^{p}\) and let f \(^{-1}\): K \(^{p}\) → K be the inverse isomorphism. Let g: K \(^{p}\) → C be the composite of the following ring homomorphisms
\end{enumerate}

\[
\mathrm{K} ^ {p} \xrightarrow {f ^ {- 1}} \mathrm{K} \xrightarrow {\varphi} \mathrm{C} / \mathrm{N} \xrightarrow {\lambda} \mathrm{C}.
\]

For every \(x \in K\), one has \(g(x^{p}) = \lambda(\varphi(x))\). Since \(\lambda(\alpha z) = \alpha^{p}\lambda(z)\) for \(\alpha \in k\) and \(z \in C/N\), the map \(g\) is \(k^{p}\)-linear. Since the extension \(K\) of \(k\) is separable, \(k(K^{p})\) is identified with \(k \otimes_{k^{p}} K^{p}\) (A, V, p.~119, remark); consequently there exists a unique homomorphism of \(k\)-algebras \(h: k(K^{p}) \to C\) which coincides with \(g\) on \(K^{p}\).

Let \((a_{i})_{i\in I}\) be a \(p\)-basis of K over \(k(\mathrm{K}^p)\) (A, V, p.~98, theorem 2); for every \(i\in I\), choose an element \(b_{i}\) of C such that \(\pi(b_{i})=\varphi(a_{i})\). One has \(h(a_{i}^{p})=g(a_{i}^{p})=\lambda(\varphi(a_{i}))=\lambda(\pi(b_{i}))=b_{i}^{p}\) for every \(i\in I\). By A, V, p.~94, remark, there exists a homomorphism of \(k\)-algebras \(\tilde{\varphi}:K\to C\), extending \(h\) and such that \(\tilde{\varphi}(a_{i})=b_{i}\) for every \(i\). One has \(\pi(\tilde{\varphi}(a_{i}))=\pi(b_{i})=\varphi(a_{i})\) for every \(i\), and \(\pi(\tilde{\varphi}(x^{p}))=\pi(h(x^{p}))=\pi(g(x^{p}))=\pi(\lambda(\varphi(x)))=\varphi(x^{p})\) for every \(x\in K\). One therefore has \(\pi\circ\tilde{\varphi}=\varphi\), which completes the proof.

COROLLARY.--- Let k be a field, K a separable extension of k, and A a linearly topologized K-algebra. If A is formally smooth over K, it is formally smooth over k.

This follows from the theorem and from prop. 3 a) of no. 2.

Remarks.- 1) Let \(k\) be a field. Every \(k\) -algebra that is étale (A, V, p.~28, def. 1) is formally smooth (loc. cit., p.~34, th. 4, d) and no. 2, prop. 3, b)).

\begin{enumerate}
\def\labelenumi{\arabic{enumi})}
\setcounter{enumi}{1}
\tightlist
\item
  We shall see below (cor. 2 of th. 2 of no. 5) that a field extension which is formally smooth is absolutely regular, hence separable (§ 6, no. 4, example 2).
\end{enumerate}

